\section{本章小结}
\label{sec:fom-summary}

本章在只有函数值和梯度的条件下回答了“朝哪走、走多远、用不用历史”。负梯度是局部最陡的方向，二次上界决定一步能走多远，$\alpha=1/L$保证每步下降。梯度下降的保证按假设分三档：非凸时梯度变小，凸时目标差按$1/T$，强凸时按$(1-1/\kappa)^T$。

病态的目标让步长被陡峭方向限制，平坦方向进展缓慢，步数与$\kappa$成正比。历史位移可以把一致的方向累积起来，把振荡的方向抵消。重球法在当前点读取梯度，在正定二次目标上把各方向的收缩因子拉平为$(\sqrt\kappa-1)/(\sqrt\kappa+1)$，但这组参数离开二次目标可能循环。Nesterov方法把梯度移到外推点，用势函数得到强凸时的$(1-1/\sqrt\kappa)^T$和凸时的$1/T^2$，适用于一般的光滑目标，但函数值不必单调，目标非凸时定理不适用。

非线性共轭梯度保留上一方向，用线搜索决定步长，不需要曲率常数。FR配合$c_2<\tfrac12$的强Wolfe条件保证下降和$\liminf\lVert\bm{g}^{(t)}\rVert_2=0$，容易卡顿；PRP能自动重置，但需要PRP+、充分下降或重启等附加规则才能谈收敛。正定二次目标上的共轭性和有限步终止，在非线性目标上都不能照搬。

比较这些方法时，保证的条件和度量必须对应，成本要按梯度读取和线搜索试探计入，迭代次数少不等于总计算少。这些方法对局部曲率的利用都是间接的：历史位移和梯度差只是曲率的影子。如果能直接计算或近似局部曲率，更新方向可以由二次模型给出，下一章由此建立Newton方向。
